% Vietnamese translation for OI Public Library
% Source-PDF: IOI中国国家候选队论文/2001/骆骥.pdf
\documentclass[11pt]{article}
\usepackage{oipl-vn}

\title{Từ ``bài toán xe buýt'' bàn về một khía cạnh của tìm kiếm sâu}
\author{Luo Ji}
\date{}

\begin{document}
\maketitle

\origtitle{On one aspect of depth-first search through the bus problem}
\sourcepath{IOI China candidate team papers / 2001 / Luo Ji.pdf}

\section{Mô tả bài toán}

Một người đứng tại một trạm xe buýt công cộng và quan sát các xe buýt đến
trạm từ 12:00 đến 12:59. Trạm này được nhiều tuyến xe buýt dùng chung; người
đó lần lượt ghi lại thời điểm xe buýt đến trạm.

\begin{itemize}
  \item Trong khoảng 12:00--12:59, các xe trên cùng một tuyến đến trạm với
  cùng một khoảng thời gian.
  \item Đơn vị thời gian là phút, đánh số từ $0$ đến $59$.
  \item Mỗi tuyến xe buýt có ít nhất hai xe đến trạm.
  \item Số tuyến xe buýt $K$ chắc chắn không vượt quá $17$, và số xe $N$ chắc
  chắn nhỏ hơn $300$.
  \item Xe từ các tuyến khác nhau có thể đến cùng một thời điểm.
  \item Hai tuyến xe buýt khác nhau có thể có cùng thời điểm xe đầu tiên đến
  trạm và cùng khoảng cách thời gian giữa các xe.
\end{itemize}

Cần lập một lịch tuyến xe buýt sao cho, với số tuyến ít nhất có thể, các thời
điểm xe buýt đến trạm thỏa mãn dữ liệu đầu vào.

Ví dụ, giả sử các xe đến như sau:
\[
\begin{array}{c|cccccccc}
\text{Số xe} & 1 & 2 & 3 & 4 & 5 & 6 & 7 & 8\\
\hline
\text{Thời điểm đến} & 0 & 3 & 5 & 13 & 14 & 14 & 21 & 25
\end{array}
\]

Khi đó có thể có một nghiệm gồm ba tuyến:
\[
\begin{array}{c|ccc|c}
\text{Tuyến 1} & 0 & 14 & & \text{khoảng cách }14\\
\text{Tuyến 2} & 3 & 14 & 25 & \text{khoảng cách }11\\
\text{Tuyến 3} & 5 & 13 & 21 & \text{khoảng cách }8
\end{array}
\]

\section{Phân tích}

Sau một loạt phân tích, ta quyết định dùng tìm kiếm sâu để giải bài này. Vì
toàn bài đều thảo luận về tìm kiếm sâu, từ đây gọi ngắn gọn là tìm kiếm.

Với bài toán này, trước hết rút ra ba yếu tố then chốt: thời gian, xe, và
tuyến. Đặc trưng của một chiếc xe là thời gian đến. Đặc trưng của một tuyến
là thời điểm xe đầu tiên và khoảng cách thời gian, tương đương với việc biết
xe thứ nhất và xe thứ hai của tuyến đó.

Từ ba yếu tố này cần xác định đối tượng tìm kiếm và chiến lược tìm kiếm. Bài
toán yêu cầu quan hệ giữa xe và tuyến, còn thời gian chủ yếu đóng vai trò mô
tả và ràng buộc. Vì vậy đối tượng tìm kiếm của bài này nên là xe hoặc tuyến.

\section{Phân tích đối tượng và chiến lược tìm kiếm}

\subsection{Đối tượng là xe}

Chiến lược sinh ra từ đối tượng này là: theo thứ tự thời gian đến trạm, lần
lượt liệt kê mỗi xe thuộc tuyến nào.

Chú ý đặc trưng của tuyến: nếu đã xác định xe thứ nhất và xe thứ hai của một
tuyến, thì tuyến đó cũng được xác định. Phác thảo tìm kiếm như sau. Theo thứ
tự thời gian đến, lần lượt xét các xe chưa xác định tuyến. Với mỗi xe, liệt kê
khả năng nó là xe đầu tiên của một tuyến mới, hoặc là xe thứ hai của một tuyến
đã có. Nếu chọn khả năng sau, ta có thể xác định tất cả các xe khác trên tuyến
đó.

Trong ví dụ đơn giản ở trên, nếu dựng cây tìm kiếm theo cách này, khi số tầng
của cây tăng lên, số nhánh mở rộng từ mỗi nút cũng dần tăng. Trực quan mà
nói, từ gốc cây tìm kiếm, số nhánh ngày càng nhiều và phần lá ngày càng rậm
rạp. Đây là đặc điểm điển hình của cây tìm kiếm khi đối tượng tìm kiếm là xe.

\subsection{Đối tượng là tuyến}

Chiến lược sinh ra từ đối tượng này là: liệt kê mỗi tuyến chứa những xe nào,
từ đó xác định tuyến.

Trên thực tế, với mỗi tuyến, chỉ cần liệt kê đặc trưng của nó: xe thứ nhất và
xe thứ hai. Đồng thời, theo tư tưởng ``chuẩn hóa thứ tự'', cố định rằng tất cả
các tuyến được sắp theo khóa là thời điểm đến của xe đầu tiên.

Phác thảo tìm kiếm như sau. Mỗi tầng của tìm kiếm xác định một tuyến: lấy xe
chưa xác định tuyến có thời điểm đến nhỏ nhất làm xe đầu tiên của tuyến mới,
rồi liệt kê xe thứ hai của tuyến này, qua đó xác định toàn bộ tuyến.

Dựa trên cùng ví dụ đơn giản, nếu dựng cây tìm kiếm của phương pháp hai, ta
thấy tính chất của nó hoàn toàn trái ngược với phương pháp một: từ gốc trở đi,
số nhánh ngày càng ít.

Hai đối tượng và chiến lược tìm kiếm này hoàn toàn khác nhau, tính chất cây
tìm kiếm cũng trái ngược. Nhìn vĩ mô, số nút trong hai cây không chênh lệch
rõ rệt. Vậy nên chọn thế nào? Lúc này cần so sánh ở mức vi mô.

\section{So sánh đối tượng và chiến lược nào tốt hơn}

Đã so sánh thì phải có tiêu chuẩn. Ở đây chọn hai tiêu chuẩn:
\begin{itemize}
  \item cách nào dễ tối ưu bằng cắt tỉa hơn;
  \item cách nào có lượng thao tác nhỏ hơn.
\end{itemize}

\subsection{So sánh khả năng tối ưu cắt tỉa}

Bài này có ba loại cắt tỉa chính.

\paragraph{Cắt tỉa tính khả thi.}
Khi đặc trưng của một tuyến đã được xác định, có thể kiểm tra tuyến đó có hợp
lệ hay không. Với phương pháp một, trong mỗi tầng, khi liệt kê xe hiện tại là
xe thứ hai của một tuyến nào đó, đều cần dùng phán đoán này. Với phương pháp
hai, mỗi tầng xác định một tuyến, cũng cần dùng phán đoán này.

Điểm mấu chốt là dựa vào dạng cây tìm kiếm: một khi phương pháp một cắt tỉa,
nó cắt bỏ được cả một mảng nhánh lớn. So với đó, hiệu quả cắt tỉa của phương
pháp hai kém hơn nhiều. Vì vậy với loại cắt tỉa này, phương pháp hai thua xa
phương pháp một.

\paragraph{Cắt tỉa bằng nghiệm tốt nhất đã biết.}
Loại cắt tỉa này phụ thuộc vào cách nào tìm được nghiệm nhanh hơn. Rõ ràng
nhờ ưu thế của cắt tỉa tính khả thi, mỗi lần cắt tỉa trong phương pháp một có
thể loại bỏ nhiều nút không khả thi, nên tốc độ tìm nghiệm không chậm hơn
phương pháp hai. Với loại cắt tỉa này, phương pháp hai cũng không tốt hơn
phương pháp một.

\paragraph{Cắt tỉa loại trùng lặp.}
Chú ý rằng yếu tố thời gian trong bài chỉ nằm trong phạm vi $0$ đến $59$, còn
số xe có thể tới $300$. Điều này cho thấy có thể có rất nhiều xe đến cùng một
thời điểm. Trong ví dụ nhỏ phía trên đã có hai xe đến lúc $14$, và việc hai xe
này thuộc tuyến nào là tương đương, nên phát sinh trùng lặp.

Với phương pháp một, nếu trong các xe đến cùng thời điểm và chưa xác định
tuyến, xe có số hiệu nhỏ hơn là xe đầu tiên của một tuyến nào đó, thì xe có số
hiệu lớn hơn cũng phải là xe đầu tiên của một tuyến. Với phương pháp hai, đối
với các xe đến cùng thời điểm và chưa xác định tuyến, ta cố định chỉ chọn xe
có số hiệu nhỏ nhất làm xe thứ hai. Hai cách đều cắt bỏ các nhánh trùng lặp,
nên hiệu quả ngang nhau.

Kết luận: nhờ tính chất tốt của cây tìm kiếm, phương pháp một có triển vọng
tối ưu cắt tỉa rộng hơn.

\subsection{So sánh lượng thao tác}

Lượng thao tác ở đây là lượng thao tác của thủ tục đệ quy chính, do vòng lặp
liệt kê của thủ tục chính và các hàm cắt tỉa quyết định.

\paragraph{Vòng lặp liệt kê của thủ tục đệ quy chính.}
Ở phương pháp một, mỗi tầng dùng vòng lặp để liệt kê một xe là xe đầu tiên của
tuyến mới hay là xe thứ hai của một tuyến đã biết. Trên cây tìm kiếm, số lượng
liệt kê của vòng lặp này phụ thuộc vào số tuyến chưa xác định xe thứ hai, lớn
nhất là $17$.

Ở phương pháp hai, mỗi tầng liệt kê xe nào là xe thứ hai. Dù đã dùng cắt tỉa
loại trùng lặp, lượng vòng lặp lớn nhất vẫn là $60$. Chỉ nhìn cận trên trực
quan đã thấy phương pháp hai không tốt hơn phương pháp một.

\paragraph{Thời gian tiêu tốn của các hàm cắt tỉa.}
Do đặc thù của bài này, các hàm cắt tỉa chính của hai phương pháp về cơ bản
không khác nhau nhiều, thời gian tiêu tốn cũng xấp xỉ nhau.

Kết luận: xét về lượng thao tác, phương pháp hai không tốt hơn phương pháp
một.

\section{Kết luận cho bài toán}

Qua phân tích vi mô ở trên, ta so sánh được hai phương pháp và rút ra kết
luận cuối cùng: phương pháp một tốt hơn phương pháp hai. Khi đã chọn được đối
tượng và chiến lược tìm kiếm, các vấn đề chính trong cài đặt như cắt tỉa cũng
đã được phân tích, việc viết chương trình trở nên tự nhiên hơn.

Thực tế cũng chứng minh kết luận này. Bảng sau so sánh thời gian chạy của hai
chương trình:
\[
\begin{array}{c|cccccc}
& \texttt{car.in1} & \texttt{car.in2} & \texttt{car.in3}
& \texttt{car.in4} & \texttt{car.in5} & \texttt{car.in6}\\
\hline
K & 3 & 5 & 8 & 10 & 15 & 17\\
\text{Phương pháp 1} & 0.01s & 0.01s & 0.05s & 0.11s & 1.48s & 1.76s\\
\text{Phương pháp 2} & 0.01s & 0.01s & 9.07s & >100s & >100s & >100s
\end{array}
\]
Có thể thấy chênh lệch ưu khuyết giữa hai cách là rất lớn.

\section{Tổng kết}

Riêng với bài này, nhìn vĩ mô rất khó biết hiệu suất của hai phương pháp khác
nhau đến mức nào. Vậy vì sao phương pháp một tốt hơn? Nguyên nhân then chốt là
nó phù hợp với tối ưu cắt tỉa trong chương trình, đồng thời lượng thao tác
nhỏ.

Vì sao lại chọn hai mặt này làm tiêu chuẩn? Ta biết rằng:
\[
  \text{thời gian tìm kiếm sâu}
  \approx
  \text{hệ số thao tác tại mỗi nút}
  \times
  \text{số nút}.
\]
Từ công thức này, rất dễ thấy ở mức vi mô rằng muốn giảm thời gian tiêu tốn
thì có hai hướng: giảm số nút, tức tối ưu cắt tỉa; và giảm hệ số thao tác tại
mỗi nút, tức giảm lượng thao tác của chương trình.

Để tăng hiệu suất tìm kiếm, ta thường lặp đi lặp lại việc cải tiến ở hai mặt
trên. Nhưng điều cũng rất quan trọng là ở tiền đề vĩ mô: làm sao để ta có thể
cắt tỉa một cách đầy đủ và hiệu quả, và làm sao để ta có thể giảm lượng thao
tác một cách đầy đủ và hiệu quả.

Do đó, việc đối tượng và chiến lược tìm kiếm tạo điều kiện tốt hay xấu cho cắt
tỉa và cho việc giảm hệ số thao tác đã trở thành tiêu chuẩn so sánh của chúng
ta.

Tuy nhiên, khi dùng hai tiêu chuẩn này, cần chú ý rằng chúng liên hệ chặt chẽ
với nhau. Muốn cắt tỉa nhiều thì thường cần hàm cắt tỉa tốt và phức tạp hơn,
nhưng điều đó lại làm tăng hệ số thao tác tại mỗi nút. Hai mặt này thống nhất
về mục tiêu, nhưng đối lập về hiệu quả cục bộ. Khi lấy chúng làm tiêu chuẩn,
cần nắm chắc cách phối hợp và tìm đúng điểm cân bằng.

Nói chung, với các bài toán tìm kiếm sâu, một đối tượng tìm kiếm và chiến
lược tìm kiếm tốt là vô cùng quan trọng. Bài viết này thông qua phân tích
``bài toán xe buýt'' đã minh họa đầy đủ điểm đó, đồng thời dựa trên công thức
thời gian tiêu tốn của tìm kiếm sâu để đề xuất tiêu chuẩn so sánh đối tượng và
chiến lược tìm kiếm: tối ưu cắt tỉa và hệ số thao tác.

Phân tích công thức thời gian của tìm kiếm sâu cũng gợi ý rằng để giải bài tốt
hơn, chỉ thay đổi ở mức vi mô là chưa đủ; trước hết cần tạo ra điều kiện vĩ
mô thuận lợi cho mục tiêu đó. Đối tượng và chiến lược tìm kiếm tốt chính là
điều kiện vĩ mô để giải thành công các bài tìm kiếm, và đó cũng là chủ ý chính
của bài viết này.

\end{document}
